\emph{(i)} A rotation $R$ about $u'$ commutes with the cyclic family shift $C$ (both are rotations about the same axis), so it preserves the eigenlines of $C$, which are $\C u'$, $\C f_{+}$ and $\C f_{-}$ (eigenvalues $1$, $\zeta^{-1}$, $\zeta$, distinct); being unitary it multiplies $f_{+}$ by a phase $e^{i\theta}$, and $Ru' = u'$. Hence $Rv(\chi) = \kappa(u' + \lambda e^{i(\chi + \theta)}f_{+}) = v(\chi + \theta)$, and as $R$ ranges over the rotations about $u'$, $\theta$ ranges over all angles. $R$ is real, so $R\bar{c} = \overline{Rc}$; it fixes $u'$ and is orthogonal, so it commutes with $P_{s} = \J/3$ and $P_{d} = 1 - P_{s}$, hence with $\sigma_{x,y} = (xP_{s} + yP_{d}) \circ \overline{\phantom{v}}$ (Proposition \claimref{synthesis.coherent-face-vector-real-slices}(i)). For the covariance, $\hat{R}a_{iA}\hat{R}^{\dagger} = \sum_{j}R_{ji}a_{jA}$, and the cofactor identity $\sum_{ij}\epsilon_{kij}R_{i'i}R_{j'j} = \sum_{m}\epsilon_{mi'j'}\,(\det R)\,(R^{-1})^{\mathsf{T}}_{mk} = \sum_{m}\epsilon_{mi'j'}R_{mk}$ (for $R \in SO(3)$) gives $\hat{R}\hat\Omega_{k}\hat{R}^{\dagger} = \sum_{m}R_{mk}\hat\Omega_{m}$ for the face operators of Theorem \claimref{families.skeleton-rank-one-plus-ordering-shift}; so $\hat{R}^{\dagger}\hat\Omega_{l}\hat{R} = \sum_{m}R_{lm}\hat\Omega_{m}$ and, the mass form being $\tfrac12\langle\hat\Omega_{k}^{\dagger}\hat\Omega_{l} + \hat\Omega_{l}\hat\Omega_{k}^{\dagger}\rangle$ in every state (Proposition \claimref{synthesis.face-mass-form-unit-plus-gram}), $H(\hat{R}\psi) = RH(\psi)R^{\mathsf{T}}$, which has the spectrum of $H(\psi)$. Also $\hat{R}\hat\Omega(c)^{\dagger}\hat{R}^{\dagger} = \hat\Omega(Rc)^{\dagger}$.

\emph{(ii)} By (i), $\hat{R}$ sends $\hat\Omega(c)^{\dagger}$, $\hat\Omega(\bar{c})^{\dagger}$ and $\hat\Omega(\sigma_{x,y}c)^{\dagger}$ to the same operators at $Rc$, fixes rotation-invariant face words, and fixes the vacuum. So $\hat{R}\psi[c] = \psi[Rc]$ for every polynomial construction, and for evolutions under a Hamiltonian built from these operators, since $\hat{R}$ conjugates that Hamiltonian at $c$ into the one at $Rc$. Then $H(\psi[Rc]) = RH(\psi[c])R^{\mathsf{T}}$ has the same spectrum, and with $c = v(\chi)$, $Rc = v(\chi + \theta)$ for every $\theta$. The coherent vector's form is that of Proposition \claimref{families.rigidity-pairing-sector-one-event}. The check builds rational rotations about $u'$ by the Cayley transform $R = (1 - A)^{-1}(1 + A)$, $A = t[u']_{\times}$, at $t = \tfrac12$ and $t = -\tfrac27$, verifies $RR^{\mathsf{T}} = 1$, $Ru' = u'$, $RC = CR$ and $R \neq 1$, and for two complex-rational source vectors $c$ and all four $\sigma_{x,y}$ verifies $\sigma_{x,y}(Rc) = R\sigma_{x,y}(c)$ and that the states $\ket{0} + \hat\Omega(c)^{\dagger}\ket{0} + \tfrac23\hat\Omega(c)^{\dagger}\hat\Omega(\sigma_{x,y}c)^{\dagger}\ket{0}$ at $c$ and at $Rc$ have normalised mass forms with identical characteristic polynomials, exactly.

\emph{(iii)} The check applies $\sum_{k}(\hat\Omega_{k}^{\dagger})^{3}$ (coefficient $\tfrac12$) to $\ket{0} + \hat\Omega(c)^{\dagger}\ket{0}$ and finds the characteristic polynomials at $c$ and $Rc$ different for both $R$ and both $c$, while at $c$ and $Cc$ they agree: the cyclic shift has $\det C = 1$, so $\hat{C}\hat\Omega_{k}^{\dagger}\hat{C}^{\dagger}$ permutes the $\hat\Omega_{k}^{\dagger}$ and the cubic sum is invariant. A rotation-invariant quadratic sum such as $\sum_{k}(\hat\Omega_{k}^{\dagger})^{2}$ would not do, being invariant under all of $SO(3)$. Since the spectrum is a continuous function of $\chi$ and a cyclic-shift-invariant construction is invariant under $\chi \mapsto \chi + 2\pi/3$, such a construction can depend on $\chi$ only modulo $2\pi/3$. That $\delta = \chi$ or $\chi - \tfrac{\pi}{2}$ in the real-slice reading is Proposition \claimref{synthesis.coherent-face-vector-real-slices}(iii), where the projected vector depends on $\delta$ directly, which is why the two readings of the masses differ.
